\cotefile{156-6}
\scanpage{66}\mspage{65}
\unsure{Ad|At} 4. \unsure{Prouvons} que c'est un $\mathrm{Inf}$ pour
$\ll$, i.e. que pour \unsure{tte} figure $\unsure{\Phi|\varphi}$, on a
\[
\Phi\ll F',\ \Phi\ll G\quad\Longrightarrow\quad\Phi\ll G'
\]

Par \unsure{At|Ax} 4, on est \unsure{ramené} au cas
où $\Phi=\deleted{\unsure{X}}Z\in\unsure{M|\mathcal{M}}$. Comme
$\deleted{X}Z\ll F'$, soit $X'$ le plus petit \unsure{élément}
de $\widetilde F^{\prime Z}=\{X'\in\widetilde{\deleted{\unsure{X}}F}^{\unsure{Y\prime|\prime}}\mid\deleted{X}Z\ll X'\}$. Il
suffit de \unsure{prouver} que $X'\deleted{\unsure{\ll}}\inserted{\leq}G'$, i.e.
$X'\ll G$. Donc on est ramené au

\uline{Lemme}\qquad Soit
\diagram{\begin{tikzcd}[ampersand replacement=\&,column sep=small,row sep=0pt]
F' \& \ll \& F\\
\rotatebox{90}{$\leq$} \& \& {}\\
X' \& \& \rotatebox{90}{$\leq$}\\
\rotatebox{90}{$\overset{\circ}{\ll}$} \& \& {}\\
Z \& \ll \& G
\end{tikzcd}}
alors \deleted{\unsure{$Z$}} $X'\ll G$.

En effet, soit $X$ l'unique \unsure{élément} de $F$ \unsure{tels|tel}
que $X'\overset{\circ}{\ll}X$. Alors par At 7 on a
\[
\deleted{X}Z\overset{\circ}{\ll}X
\]
donc $X$ est le plus petit \unsure{élément} de
$\widetilde F^{Z}$. Comme $Z\ll G$, $\exists\,Y\in\widetilde G$ avec
$\deleted{X}Z\ll Y$, \inserted{(par At 4)} \unsure{donc} \deleted{\unsure{$X$}} \unsure{comme} $\widetilde G\subset\unsure{\widehat F|\widetilde F}$ d'où $Y\in\widetilde F$
on a par \unsure{définition} de $X$, $X\leq Y$, \deleted{\unsure{d'où} $X\in\widetilde G_{\unsure{z|2|Z}}$}
donc $\deleted{Z\ll X\unsure{<|\leq}}$ \unsure{ainsi} $X'\ll X\leq Y\leq G$
d'où $X'\ll G$, qed.

\scanpage{67}\mspage{66}
\noindent\llap{\smash{\rule[-4.8\baselineskip]{0.4pt}{4.2\baselineskip}}\enspace}\uline{Corollaire} 1\quad \deleted{L'application \unsure{$G$} \unsure{a aussi}}
\inserted{Si $F'$ est aussi une sous}
figure i.e. $F'\leq F$, alors $G'$ est une
\[
\deleted{G'=\mathrm{Inf}^{\leq}(F',G)}\tag*{\deleted{(1.61)}}
\]
sous-figure de $G$, \unsure{donc} de $F$, et
c'est $F'\cap G$ (sous-figure \unsure{intersection})

\unsure{Formel}, en \unsure{renversant} les rôles de
$F'$ et $G$ \unsure{--} et le fait que $G'=F'\cap G$
se voit sur (1.60).

% slanted
\margin{Notation $F'.G$}
\uline{Corollaire 2}\quad \inserted{Si $F'\ll F$,} L'application $G\longmapsto G'=F'.G$
\[
\mathrm{Sous\,fig}\,(F)\longrightarrow\mathrm{Sous\,fig}\,(F')\tag*{\unsure{(1.51)|(1.61)}}
\]
\unsure{Commute|commute} aux $\mathrm{Sup}$ \deleted{\unsure{qui|qu}} et aux $\mathrm{Inf}$.

\deleted{En} En effet, en \unsure{termes} de l'application
canonique \unsure{croissante}
\[
\unsure{\varphi|\phi}:\widetilde{\unsure{F|\mathbb F}}'\longrightarrow\widetilde F
\]
d'ens \unsure{ordonnés}, cette application s'identifie
à l'application \unsure{correspondante} \unsure{au|en} sens inverse
\[
\begin{aligned}
\varphi^{*}:{}&\unsure{\mathcal{P}f|\mathcal{P}_{f}}(\widetilde F)\longrightarrow\mathcal{P}f(\widetilde F')\\
&\unsure{A|\mathbb A}\longmapsto\varphi^{-1}(A)
\end{aligned}
\quad,
\]
qui a les propriétés en question.

\uline{Corollaire 3}\quad Soit $F'\ll F$, et soit $(F_i)$ une
famille de sous-figures de $F$ telles que
\margin{Notation $\bigvee F_i$}
$F=\mathrm{Sup}\,F_i$. Soit $F_{ij}=F_i\cap F_j$. \unsure{Considérons} les
\scanpage{68}\mspage{67}
sous figures $F'_i$, $F'_{ij}$ de $F'$ \unsure{induites} par
les $F_i$, $F_{ij}$. Alors on a
\[
\deleted{F'=\mathrm{Sup}^{\leq}F'_i,\qquad F'_i\cap F'_j=F'_{ij}}.
\]

(i)\quad Chaque $F'_i$ est un \deleted{\unsure{des|de}}raffinement de $F_i$.

(ii)\quad Le raffinement \unsure{induit} par $F'_i$ sur $F_{ij}$\inserted{$=F_i\cap F_j$} est $F'_{ij}$
(\unsure{donc} égal à celui induit par $F'_j$ sur $F_{ij}$)
% slanted
\margin{\tikz[baseline]{\draw[->] (0.16,.52) to[bend right=8] (0,0);}\quad\uline{NB} Si pour \unsure{tt} $i,j$, on se donne \deleted{\unsure{des} \unsure{recou|recouv}} $F_{ij\unsure{h|k}}\leq F_{ij}$ tels que $\mathrm{Sup}\,F_{ijh}=F_{ij}$ alors (ii) équivaut à $F'_i\mid F_{ijh}=F'_j\mid F_{\unsure{jh|ijh}}$.}

(iii)\quad $F'=\mathrm{Sup}\,F'_i$, \deleted{\unsure{et}} \unsure{est|et} c'est un
raffinement de $F$.

On a donc un \uline{principe de recollement}
pour \unsure{des|les} raffinements de $F$\,: un tel
raffinement est \unsure{connu} (par (iii)) quand
on se donne une ``famille compatible''
de \unsure{raffinements} des $F_i$. La question est
si on peut se donner celle-ci \unsure{arbitrairement}.
C'est ce qui va être \unsure{assuré}
par l'axiome suivant

\uline{\unsure{At|Ax} 1\deleted{\unsure{0}}\unsure{1|0}}\quad (\uline{Axiome de recollement} des
raffinements) Soit $F\in\unsure{\mathcal F|\mathscr F}$, et pour \unsure{tt}
$X\in\widetilde F$, soit $F'_X\unsure{\ll\!-\!|\twoheadleftarrow}X$ \deleted{\unsure{réf|raf}} un raffinement
de $X$. On \unsure{suppose} que pour $X,Y\in\widetilde F$,
\deleted{$X\unsure{\neq|\pm}Y$, on ait $X\unsure{\leq|<}Y$}\inserted{$Y<X$}, on ait
\[
\deleted{F'_X\mid X\cap Y=}\qquad F'_Y=F'_X\mid Y
\]
Alors $\exists\,F'\ll F$, tel que pour tout $X\in\widetilde F$, $F'_X=F'\mid X$.

\scanpage{69}\mspage{68}
\uline{NB}\quad Par le corollaire \unsure{précédent}, cet $F'$ est
unique. On l'appelle le raffinement de $F$
\uline{obtenu par recollement} à partir des raffinements
$F_X$ de ses \unsure{multistrates}.

\uline{NB}\quad \deleted{Cet axiome n'est pas \unsure{bien}}
\inserted{Le \deleted{\unsure{de|du}} contenu \unsure{essentiel} de}
cet axiome, c'est que les $F'_X$ soient
compatibles entre eux\inserted{deux à deux}, de façon à pouvoir
former $\mathop{\mathrm{Sup}}\limits_{X\in\widetilde F}F'_X=F'$ (du moins, on \unsure{pourra|pourrait} le
faire, \unsure{indépendamment} de \unsure{At|Ax} 11, si on suppose
$\widetilde F$ fini ...). La difficulté est de comparer
$F'_X$ et $F'_Y$ quand on n'a pas $\deleted{\unsure{X}}Y\leq X$ ou $X\leq Y$.

Donc le \unsure{noeud|nœud} de At 10 est ceci

\noindent\llap{\smash{\rule[-5.1\baselineskip]{0.4pt}{5.2\baselineskip}}\enspace}\uline{\unsure{Cor|Cor.} \unsure{1|\textbar}}\inserted{(\uline{At} 1\deleted{\unsure{0}}\unsure{1|0}$'$)}\quad Soient $X,Y\in\unsure{M|\mathcal{M}}$ avec $X\mathrel{\rlap{$<$}{>}}Y$, soit
\deleted{$\unsure{F|E}=X\vee Y$,} $\deleted{\unsure{Z}}\deleted{\unsure{G}}T=X\cap Y$. Soit $\deleted{F}X'\ll X$, $\deleted{G}Y'\ll Y$
\unsure{tels} que $\deleted{F}X'\mid T=\deleted{G}Y'\mid T$ (ou, ce qui
revient au $\unsure{\overline{\mathrm m}}$, \unsure{pour} tout $Z\in\widetilde X\cap\widetilde Y$,
$\deleted{F}X'\mid\unsure{2|Z}=\deleted{G}Y'\mid Z$). Alors $\deleted{F}X'\mathrel{\rlap{$<$}{>}}\deleted{G}Y'$.

Ceci est une variante à peine
\unsure{affaiblie} de At 1\deleted{\unsure{0}}1$^{\unsure{\prime|\checkmark}}$\unsure{.} Elle implique At 10
au moins pour $\widetilde F$ fini --- plus \unsure{généralt|généralement},
elle implique \unsure{de|dans} At 1\deleted{\unsure{0}}1 que les $F'_X$ sont
\unsure{compatibles|admissibles} entre eux, et \uline{si} ce $\mathrm{Sup}\,F'_X$ existe
\scanpage{70}\mspage{69}
(ce qui est \unsure{automatique} si on suppose $\unsure{\widehat F|\widetilde F}$ fini),
alors je dis que le \deleted{\unsure{$F'_X$}} $\mathrm{Sup}$ $F'$ répond
à la question. En effet, on aura \unsure{$\Gamma$|$F$}
\[
\widetilde F'=\deleted{\unsure{S}}\bigcup_X\widetilde F'_X\quad,
\]
et on aura $F'\ll\unsure{F|\mathbb F|\overline F}$ par \unsure{At|Ax} 4. Pour
prouver que $F'\mid X=F'_X$\inserted{$\forall X\in\widetilde F$},
\deleted{\unsure{on note que}}
\[
\deleted{F'_X=\mathop{\mathrm{Sup}}\limits_{Z\in\unsure{F|\widetilde F}}\bigl(F'_Z\mid X\bigr)}
\]
% slanted
\margin{\deleted{\unsure{réduction} des raffinements \inserted{$F'_Z$} de $F$ sur \unsure{sous}-figure $X$ de $F$}}
\deleted{Et par la condition \unsure{mise|mises} sur le \unsure{syst.} des $F'_T$, on aura}
\diagram{\begin{tikzcd}[ampersand replacement=\&,column sep=2pt,row sep=1pt,cells={nodes={font=\scriptsize,inner sep=1pt}}]
{} \& {} \& {} \& {} \& {} \& {} \& \deleted{F'} \& \deleted{\ll} \& \deleted{F} \& {} \& {}\\
{} \& {} \& {} \& {} \& {} \& {} \& \deleted{\rotatebox{90}{$\leq$}} \& \deleted{\unsure{F'_Z|F_Z}} \& \deleted{\rotatebox{80}{$\leq$}} \& \deleted{\rotatebox{105}{$\leq$}} \& \deleted{F'_Z\mid X}\\
\deleted{\mathop{\mathrm{Sup}}\limits_Z F'_Z} \arrow[dd,dash,bend right=16] \& \deleted{=} \arrow[dd,dash,bend left=18] \& \deleted{F'} \& \deleted{\ll} \& \deleted{F} \& {} \& \deleted{F'_Z} \& \deleted{\ll} \& \deleted{Z} \& \deleted{X} \& {}\\
\deleted{\unsure{\rotatebox{90}{$\leq$}}} \& {} \& \deleted{\rotatebox{90}{$\leq$}} \& {} \& \deleted{\rotatebox{90}{$\leq$}} \& {} \& {} \& {} \& \deleted{\rotatebox{125}{$\leq$}} \& \deleted{\rotatebox{70}{$\leq$}} \& {}\\
\deleted{F'_Z\mid X} \& \deleted{\leq} \& \deleted{F'\unsure{\mathord{.}|\cdot}\mid X} \& \deleted{\ll} \& \deleted{X} \& {} \& {} \& {} \& {} \& |[alias=inkIntersection]| \deleted{X\cap Z} \& {}
\arrow[from=1-5,to=3-5,dash,start anchor=north,end anchor=center]
\arrow[from=1-8,to=3-9,dash,start anchor=north,end anchor=south west,bend right=7]
\arrow[from=1-10,to=3-10,dash,start anchor=north east,end anchor=east,bend right=5]
\arrow[from=3-3,to=5-1,dash,bend right=12]
\arrow[from=3-4,to=5-5,dash,bend right=25]
\arrow[from=3-5,to=2-11,dash,start anchor=center,end anchor=east,
  to path={.. controls +(0,-8mm) and +(-9mm,4mm) .. (inkIntersection.south west)
  .. controls +(6mm,-2mm) and +(2mm,5mm) .. (inkIntersection.east)
  .. controls +(-2mm,-2mm) and +(-4mm,-7mm) .. (\tikztotarget) \tikztonodes}]
\end{tikzcd}}

On regarde la \unsure{situation ensembliste}
\diagram{\begin{tikzcd}[ampersand replacement=\&,row sep=0pt,column sep=large]
\widetilde F'=\bigcup\widetilde F'_X \arrow[r] \& \widetilde F\,\unsure{-}\\
{} \& \rotatebox{90}{$\subset$}\\
{} \& \widetilde X
\end{tikzcd}}
où $\widetilde F$ est réunion des $\widetilde X$, et où $\widetilde F'$ a été
\unsure{obtenu} par recollement \deleted{\unsure{des|du}} \inserted{du syst. des} \unsure{ens} $\widetilde F'_X$
sur la famille admissible des $\widetilde X$ \unsure{recouvrant}
$\widetilde F$ ...

\uline{Corollaire} \uline{2}\quad Sous les conditions du Cor. 3
de la page 66, si on se donne un
\uline{\unsure{syst} compatible} $(F'_i)$ de raffinements des
$F_i$ (définir ce terme), alors
\scanpage{71}\mspage{70}
il existe un \unsure{raffinement} $F'$ de $F$ (\unsure{nec.|néc.} unique)
qui l'induit.

Dans le cas où la famille \deleted{\unsure{de|d}}d'indices
$I$ est \uline{fini}, ceci résulte déjà du seul
axiome affaibli \uline{\unsure{At|Ax}} $1\deleted{0}1'$ \unsure{.|--}

\unsure{Cas|En} particulier

\uline{Cor. 3}\quad Soient $F,G$ deux figures avec
$F\parallel G$, et considérons la figure somme\inserted{$H=F\amalg G$}.
Alors

(i) Si $F'\ll F$, $G'\ll G$, on a $F'\parallel G'$, et
\[
H'\overset{\unsure{\mathrm{def}|\text{déf}|\mathrm{dfn}}}{=}F'\amalg G'\deleted{\unsure{<}}\ll H=F\amalg G.
\]

\inserted{(ii)} On \unsure{trouve} ainsi une bijection
\[
\begin{aligned}
\mathrm{Raff}(F)\times\mathrm{Raff}(G)&\xrightarrow{\sim}\mathrm{Raff}(F\amalg G)\\
(F',G')&\unsure{\mathrel{\tikz[baseline=-.5ex]{\draw (0,-.55ex)--(0,.55ex); \draw (0,0) .. controls (.8em,.02ex) and (1.45em,-.24ex) .. (2em,.12ex) -- (2em,-.24ex);}}|\longmapsto}H'=F'\amalg G'
\end{aligned}
\]
la bijection en sens \unsure{inverse} est \unsure{donnée} par
\[
H'\longmapsto (H'\mid F,H'\mid G) \qquad .
\]

Se \unsure{généralise} au cas d'une famille
$F$, \uline{somme} \uline{disjointe} de sous-familles $F_i$ \unsure{--|.}

On \unsure{retrouve} en particulier la
\deleted{\unsure{lo|re}}\unsure{relation} $(1.58)$ p. 59
\[
F\deleted{\unsure{<}}\parallel G,\qquad F'\ll F,\ G'\ll G
\quad\Longrightarrow\quad F'\parallel G'.
\]

\scanpage{72}\mspage{71}
\deleted{\uline{\unsure{Cor. 4}}}

% vertical pen rule at the left of the definition
\noindent\llap{\smash{\rule[-1.7\baselineskip]{0.4pt}{1.6\baselineskip}}\enspace}\uline{\unsure{Definition|Définition}}\quad Une figure \deleted{\unsure{est} $F$} $\unsure{\Phi|\varphi}$ est dite
\uline{discrète} si toutes ses \deleted{\unsure{mult|multi}} strates sont
des lieux, i.e. $\widetilde\Phi\subset\unsure{\mathcal L|L}$.

\deleted{\uline{\unsure{Cor 4|Cor. 4}}}\quad \deleted{\unsure{Elle}}Une telle figure est connue
par la \unsure{connaissance} de la partie $\widetilde\Phi=A$
de $\mathcal L$. Celle-ci est \unsure{formée} de
lieux deux à deux disjoints, et \deleted{\unsure{pour}}
pour que une partie \uline{finie} de $\mathcal L$ ``soit''
une figure, il (f. et) s. qu'\unsure{elle|elles}
\deleted{\unsure{points|point} \unsure{soi|si|qui} \illegible{1} \unsure{soient}}
\inserted{soit \unsure{formée} de} lieux deux à deux
disjoints\unsure{..|.} \deleted{Plus \unsure{gén.|gén} , \unsure{sans|sa} condition de \unsure{finitude}.}
\inserted{\unsure{Notons|Notant|Noter} que la relation
d'ordre $\leq$ ou $\ll$ sur $\unsure{\widetilde\Phi|\overline{\widetilde\Phi}}=A$ est la relation}
discrète (c'est ainsi sur $\mathcal L$ !) donc
les parties fermées sont les parties
quelconques, donc les \unsure{s/figures|s.figures} \deleted{sont}
\unsure{correspondent} aux parties quelconques.

Plus \unsure{gén|gén.}\textsuperscript{t}, si $F$ est une figure,
on s'intéresse aux sous figures discrètes.
Elles sont décrites par \unsure{certaines} \deleted{$A\subset$}
parties $A$ de $\mathcal L$.

\scanpage{73}\mspage{72}
\uline{Cor 4}\quad \inserted{Soit $F$ une figure\unsure{,|.}} Pour qu'une partie $A$ de
$\unsure{\mathcal L|L}$ \unsure{définisse} une figure discrète\inserted{$\unsure{\Phi|\varphi}$} et qu'on
ait $\Phi\ll F$, il \unsure{f. et. s.|f. et s.} que
\[
\begin{aligned}
\text{(i)}\quad &\Phi\subset\mathrm{omb}(F)\\
\text{(ii)}\quad &\forall X\in\widetilde F,\qquad \Phi\cap\mathrm{omb}(X)\quad \text{``soit une figure''}
\end{aligned}
\]
(p.ex soit fini et \unsure{formé} de \deleted{\unsure{élém.|él}}\deleted{\unsure{de|d'}} \unsure{lieux} deux à deux
disjoints).

Ceci implique \unsure{notamment} que à la
page \unsure{67|62}, dans l'énoncé de \unsure{At|Ax} $9'$
(critère suffisant pour que $X\mathrel{\rlap{$<$}{>}}Y$), la
condition donnée \deleted{soit} en termes
de disjonction de lieux, est aussi
\uline{nécessaire} [mais cela résulte aussi directement
de At 10]

% slanted
\margin{At 10 \unsure{posé} après, \unsure{page} 64.}
% two left pen rules: the outer one spans statement and proof
\noindent\llap{\smash{\rule[-9.2\baselineskip]{0.4pt}{9.8\baselineskip}}\quad}%
\uline{Corollaire $\deleted{\unsure{4}}5$}\quad Soient $F$ une figure, $G$ une
\llap{\smash{\rule[-1.4\baselineskip]{0.4pt}{1.8\baselineskip}}\enspace}sous figure $\neq F$. Alors $\exists x\in\mathcal L$ tel que
$x\unsure{\ll|\underline{\ll}}\unsure{\mathbb F|F}$ (i.e. $x\in\mathrm{omb}(\mathbb F)$) et $x\parallel G$.

Soit $X\in\widetilde{\mathbb F}\setminus\widetilde G$, et $x\in\mathrm{omb}(X)^\circ$ \unsure{(148)|(At 8)|(A48)|(1.48)}
\unsure{i.e.|i.e.e} $x\in\mathcal L$, $x\unsure{\overset{\circ}{\ll}|\overset{\circ}{\underline{\ll}}}X$. Je dis qu'il fait l'affaire.
Pour voir qu'il est disjoint de $G$, on peut
soit invoquer At 9 (critère de disjonction pour les
lieux) et At 10, (critère de disjonction pour les lieux),
soit se \unsure{référer} à \deleted{\unsure{un|ce} \unsure{même|axiome}} l'axiome de recollement, pour
\unsure{construire} $G\amalg x$ comme \unsure{raffinement de|raffinant} $F$ ...
